\hitFrontHeading{摘\quad 要}

高效能源系统的规划、运行与维护正在经历由经验驱动向数据驱动转变的过程。面向综合能源站、热电联供系统和分布式储能装置，现场传感数据、运行日志、天气预报和设备台账之间存在采样频率不同、噪声水平不同、语义粒度不同等问题。本文以公开场景构造博士论文示例，围绕多源数据融合、状态建模、调度优化和结果解释四个环节展开，展示 `thesis-hit-doctor` 模板在哈尔滨工业大学博士学位论文结构中的完整排版能力。

本文首先建立多源数据质量分层框架，将设备状态、环境扰动和运行约束统一映射到可计算特征空间；其次构建能效预测模型，用于描述不同负荷条件下系统运行效率的变化趋势；再次设计滚动调度方法，在满足设备安全边界的前提下平衡能源消耗、峰值负荷和运行稳定性；最后给出可解释性分析流程，使模型结果能够服务于工程复核和论文写作。示例正文采用可公开分享的占位数据和演示表图，不涉及真实企业、个人或保密项目。

排版层面，模板覆盖中文封面、中文题名页、英文题名页、中英文摘要、目录、英文目录、章节正文、图表、公式、结论、参考文献、创新成果、评阅与答辩信息、原创性声明、致谢和个人简历等页面。各部分均通过 `bensz-thesis` 统一构建链路生成 PDF，便于后续按 HIT 官方 Word 范例继续做像素级细化。

\hitKeywordsZhLine
